\documentclass[../../main.tex]{subfiles}
\begin{document}
% SOURCE p=179; printed=164; object=ch05-baire-exercises-title; region=top
\subsection*{习题}
% SOURCE p=179; printed=164; object=ch05-ex27-meager-full-measure; number=27; region=top
\begin{exercise}\label{exercise:ch05-meager-full-measure}
存在 $\mathbb R$ 的疏子集，其补集的 Lebesgue 测度为零。
\end{exercise}
% SUPPLEMENT: id=v14-ch05-ex27-solution; parent=ch05-ex27-meager-full-measure; kind=solution+motivation+hypotheses; origin=independent-derivation
\begin{solution}
\textbf{补写解答（重构新增）。}列举 $\mathbb Q=\{q_k:k\ge1\}$。对每个 $j\ge1$，令
\begin{align*}
G_j=\bigcup_{k=1}^{\infty}\bigl(q_k-2^{-j-k-1},q_k+2^{-j-k-1}\bigr).
\end{align*}
它开且包含 $\mathbb Q$，所以稠密；由可数次可加性，$m(G_j)\le\sum_{k\ge1}2^{-j-k}=2^{-j}$。令 $G=\bigcap_{j\ge1}G_j$，它是 Borel 集，且每个 $j$ 都给 $0\le m(G)\le2^{-j}$，故 $m(G)=0$。令 $E=\mathbb R\setminus G=\bigcup_{j\ge1}(\mathbb R\setminus G_j)$。各补集闭且内部为空，所以无处稠密；故 $E$ 疏，而其补集 $G$ 为零测集，即为所求。Baire 定理还给 $G$ 稠密，说明一个集合可以既剩余又零测。
\end{solution}
\begin{remark}
\textbf{动机与条件（重构新增）。}围住所有有理点的小区间保证开稠密，同时把总长度压至任意小。取交使测度降到零，取补使集合成为可数个无处稠密闭集的并。这个例子把拓扑小性与测度小性明确地区分开来；两者不能互相替代。
\end{remark}
% END SUPPLEMENT: id=v14-ch05-ex27-solution
% SOURCE p=179; printed=164; object=ch05-ex28-lch-baire; number=28; region=upper
\begin{exercise}\label{exercise:ch05-lch-baire}
若将 Baire 纲定理中“$X$ 为完备度量空间”改成“$X$ 为 LCH 空间”，定理仍成立。（证明类似；\cref{proposition:ch04-compact-fip}代替完备性。）
\end{exercise}
% SUPPLEMENT: id=v14-ch05-ex28-solution; parent=ch05-ex28-lch-baire; kind=solution+motivation+hypotheses; origin=independent-derivation
\begin{solution}
\textbf{补写解答（重构新增）。}先证开稠密集列 $U_n$ 的交稠密。给定非空开集 $W$，$W\cap U_1$ 非空。由 LCH 空间的小紧邻域性质，选非空紧邻域 $K_1\subseteq W\cap U_1$。递归地，$K_n^\circ\cap U_{n+1}$ 非空且开，故可在其中选一个非空紧邻域 $K_{n+1}$。于是
\begin{align*}
\varnothing\ne K_{n+1}\subseteq K_n^\circ\cap U_{n+1}\subseteq K_n.
\end{align*}
Hausdorff 性使各紧集 $K_n$ 闭，因而它们在紧空间 $K_1$ 中是非空嵌套闭集，具有有限交性质。由\cref{proposition:ch04-compact-fip}，其交非空。任取交中的 $x$，则 $x\in W$ 且对每个 $n$ 都有 $x\in U_n$，证明稠密。

若再假设 $X\ne\varnothing$，对任意无处稠密集列 $E_n$ 应用上述结论于 $X\setminus\overline E_n$，得到交非空，故 $\bigcup_nE_n\ne X$。与原 Baire 定理一样，第二问需要非空条件；第一问在空空间仍成立。
\end{solution}
\begin{remark}
\textbf{动机与逐条件说明（重构新增）。}局部紧性提供有内部的紧邻域，Hausdorff 性保证这些紧集闭，并支持小紧邻域的缩小。有限交性质保证嵌套构造留下一个点，代替度量证明中“半径趋零加完备性”的步骤。本证明不需度量，也不声称交点唯一。
\end{remark}
% END SUPPLEMENT: id=v14-ch05-ex28-solution
% SOURCE p=179; printed=164; object=ch05-ex29-weighted-incomplete; number=29; region=upper-middle
\begin{exercise}\label{exercise:ch05-weighted-incomplete}
\begin{definition*}[本题的序列空间]\phantomsection\label{definition:ch05-weighted-sequence-domain}
令 $\mathcal Y=L^1(\mu)$，其中 $\mu$ 为 $\mathbb N$ 上的计数测度；令 $\mathcal X=\{f\in\mathcal Y:\sum_{n=1}^{\infty}n|f(n)|<\infty\}$，并赋予 $L^1$ 范数。
\end{definition*}
\begin{enumerate}
\item\label{item:ch05-weighted-dense} $\mathcal X$ 是 $\mathcal Y$ 的真稠密子空间，因此 $\mathcal X$ 不完备。
\item\label{item:ch05-weighted-closed-unbounded} 定义 $T:\mathcal X\to\mathcal Y$ 为 $Tf(n)=nf(n)$。则 $T$ 闭而无界。
\item\label{item:ch05-weighted-inverse-not-open} 令 $S=T^{-1}$。则 $S:\mathcal Y\to\mathcal X$ 有界且满射，但不开。
\end{enumerate}
\end{exercise}
% SUPPLEMENT: id=v14-ch05-ex29-solution; parent=ch05-ex29-weighted-incomplete; kind=solution+well-definedness+motivation+hypotheses; origin=independent-derivation
\begin{solution}
\textbf{逐项补写（重构新增）。}
\begin{enumerate}
\item 对应\subcref{exercise:ch05-weighted-incomplete}{item:ch05-weighted-dense}：加权绝对和对线性组合仍有限，所以 $\mathcal X$ 是子空间。它包含所有有限支集序列，这些序列由截断在 $\ell^1$ 中稠密。而 $f(n)=n^{-2}$ 属于 $\mathcal Y$，却因 $\sum_nnf(n)=\sum_nn^{-1}=\infty$ 不属于 $\mathcal X$，所以为真子空间。该 $f$ 的有限截断在 $\mathcal X$ 中为 Cauchy 列，却在 $\mathcal Y$ 中趋于不在 $\mathcal X$ 的 $f$；极限唯一性说明在 $\mathcal X$ 中没有极限，故不完备。
\item 对应\subcref{exercise:ch05-weighted-incomplete}{item:ch05-weighted-closed-unbounded}：定义 $\mathcal X$ 的加权条件正好保证 $Tf\in\mathcal Y$，且 $T$ 线性。若 $f_j\to f$ 于 $\mathcal X$ 且 $Tf_j\to g$ 于 $\mathcal Y$，$\ell^1$ 收敛蕴含每个坐标收敛，所以 $g(n)=\lim_jnf_j(n)=nf(n)$，即 $g=Tf$。故图像闭。对第 $n$ 个单位坐标序列 $e_n$，$\|e_n\|_1=1$ 而 $\|Te_n\|_1=n$，因此无界。
\item 对应\subcref{exercise:ch05-weighted-incomplete}{item:ch05-weighted-inverse-not-open}：对 $g\in\mathcal Y$，定义 $Sg(n)=g(n)/n$，则 $\sum_nn|Sg(n)|=\sum_n|g(n)|<\infty$，所以 $Sg\in\mathcal X$；并且 $TSg=g$、$STf=f$。故确为双射逆映射，且 $\|Sg\|_1\le\|g\|_1$。若 $S$ 开，其逆 $T$ 连续，进而有界，与前问矛盾。因此 $S$ 不开。
\end{enumerate}
\end{solution}
\begin{remark}
\textbf{动机与条件（重构新增）。}加权条件限制定义域，却仍用未加权范数量大小，使高位单位坐标被 $T$ 任意放大。$\mathcal Y=\ell^1$ 完备，失败的正是 $T$ 的定义域或 $S$ 的值域完备性。若改用加权范数，空间性质和算子有界性都会改变，不能与本题混用。
\end{remark}
% END SUPPLEMENT: id=v14-ch05-ex29-solution
% SOURCE p=179; printed=164; object=ch05-ex30-derivative-incomplete; number=30; region=middle
\begin{exercise}\label{exercise:ch05-derivative-incomplete}
令 $\mathcal Y=C([0,1])$、$\mathcal X=C^1([0,1])$，两者均赋一致范数。
\begin{enumerate}
\item\label{item:ch05-c1-uniform-incomplete} $\mathcal X$ 不完备。
\item\label{item:ch05-derivative-closed-unbounded} 映射 $\mathrm d/\mathrm dx:\mathcal X\to\mathcal Y$ 闭（见\cref{exercise:ch05-ck-banach}），但无界。
\end{enumerate}
\end{exercise}
% SUPPLEMENT: id=v14-ch05-ex30-solution; parent=ch05-ex30-derivative-incomplete; kind=solution+motivation+hypotheses; origin=independent-derivation
\begin{solution}
\textbf{逐项补写（重构新增）。}
\begin{enumerate}
\item 对应\subcref{exercise:ch05-derivative-incomplete}{item:ch05-c1-uniform-incomplete}：令 $f_j(t)=\sqrt{(t-1/2)^2+j^{-1}}$，它属于 $C^1([0,1])$，且 $0\le f_j(t)-|t-1/2|\le j^{-1/2}$。所以 $f_j$ 一致趋于 $|t-1/2|$，后者在 $1/2$ 不可微。因而 $f_j$ 在一致范数下为 Cauchy 列，却无 $\mathcal X$ 内的极限，证明不完备。
\item 对应\subcref{exercise:ch05-derivative-incomplete}{item:ch05-derivative-closed-unbounded}：若 $f_j\to f$ 于 $\mathcal X$ 的一致范数且 $f_j'\to g$ 于 $\mathcal Y$，由
\begin{align*}
f_j(t)-f_j(0)=\int_0^t f_j'(s)\,\mathrm ds
\end{align*}
并对两边取极限，得到 $f(t)-f(0)=\int_0^tg(s)\,\mathrm ds$。因 $g$ 连续，微积分基本定理给 $f'=g$，故图像闭。取 $f_n(t)=t^n$，其一致范数为1，而导数一致范数为 $n$，所以微分算子无界。
\end{enumerate}
\end{solution}
\begin{remark}
\textbf{动机与条件（重构新增）。}一致控制函数值不控制变化速度，故微分可能任意放大误差；但若函数和导数都已有一致极限，积分恒等式仍保持导数关系，这正是图像闭性。与\cref{exercise:ch05-ck-banach}中的 $C^1$ 范数不同，本题没有把导数范数加进去，因而定义域不完备。
\end{remark}
% END SUPPLEMENT: id=v14-ch05-ex30-solution
% SOURCE p=179; printed=164; object=ch05-ex31-unbounded-graph-space; number=31; region=middle
\begin{exercise}\label{exercise:ch05-unbounded-graph-space}
设 $\mathcal X,\mathcal Y$ 为 Banach 空间，$S:\mathcal X\to\mathcal Y$ 为无界线性映射（其存在性见\cref{section:ch05-notes}）。令 $\Gamma(S)$ 为 $S$ 的图像，它是 $\mathcal X\times\mathcal Y$ 的子空间。
\begin{enumerate}
\item\label{item:ch05-unbounded-graph-incomplete} $\Gamma(S)$ 不完备。
\item\label{item:ch05-graph-parametrization-closed} 定义 $T:\mathcal X\to\Gamma(S)$ 为 $Tx=(x,Sx)$。则 $T$ 闭而无界。
\item\label{item:ch05-graph-projection-not-open} $T^{-1}:\Gamma(S)\to\mathcal X$ 有界且满射，但不开。
\end{enumerate}
\end{exercise}
% SUPPLEMENT: id=v14-ch05-ex31-solution; parent=ch05-ex31-unbounded-graph-space; kind=solution+motivation+hypotheses; origin=independent-derivation
\begin{solution}
\textbf{逐项补写（重构新增）。}乘积及其子空间均取最大范数。
\begin{enumerate}
\item 对应\subcref{exercise:ch05-unbounded-graph-space}{item:ch05-unbounded-graph-incomplete}：若图像完备，它作为度量空间的完备子空间必闭：对收敛到环境点的图像序列，完备性给图像内极限，唯一性使两极限相同。于是 $S$ 图像闭，由闭图像定理得到 $S$ 有界，矛盾。因此图像不完备。
\item 对应\subcref{exercise:ch05-unbounded-graph-space}{item:ch05-graph-parametrization-closed}：若 $x_n\to x$ 于 $\mathcal X$ 且 $Tx_n\to z$ 于 $\Gamma(S)$，则 $z=(u,Su)$ 对某 $u\in\mathcal X$ 成立。第一坐标收敛给 $x_n\to u$，故 $u=x$，所以 $z=(x,Sx)=Tx$，证明 $T$ 闭。又 $\|Tx\|=\max(\|x\|,\|Sx\|)\ge\|Sx\|$，所以 $T$ 不可能有界。
\item 对应\subcref{exercise:ch05-unbounded-graph-space}{item:ch05-graph-projection-not-open}：$T$ 显然双射，逆为第一坐标投影，满足 $\|T^{-1}(x,Sx)\|=\|x\|\le\|(x,Sx)\|$，故有界且满射。若它开，逆 $T$ 连续，进而有界，再与前问矛盾。
\end{enumerate}
\end{solution}
\begin{remark}
\textbf{动机、条件与存在性补充（重构新增）。}这里定义域 $\mathcal X$ 完备，失败的是目标图像空间的完备性，因而补足前两题所显示的另一端反例。$S$ 在整个 $\mathcal X$ 上线性，使图像参数化为双射；$S$ 无界使图像不闭。

这样的 $S$ 的确存在。例如在 $\ell^1$ 中，将线性无关族 $\{e_n:n\ge1\}$ 延拓为 Hamel 基：包含它的线性无关族按包含排序，链并仍线性无关，Zorn 引理给极大族，极大性使它张成全空间。在该基上规定 $S(e_n)=n$，其他基向量的值为0，再按每个向量唯一的有限基展开定义线性泛函 $S:\ell^1\to\mathbb C$（实域同理）。因 $\|e_n\|_1=1$ 而 $|S(e_n)|=n$，它无界，且两端均为 Banach 空间。这一构造已独立证明原文预告的存在性。
\end{remark}
% END SUPPLEMENT: id=v14-ch05-ex31-solution
% SOURCE p=179; printed=164; object=ch05-ex32-comparable-banach-norms; number=32; region=middle-lower
\begin{exercise}\label{exercise:ch05-comparable-banach-norms}
设 $\|\cdot\|_1,\|\cdot\|_2$ 为向量空间 $\mathcal X$ 上的范数，且 $\|\cdot\|_1\le\|\cdot\|_2$。若 $\mathcal X$ 关于两个范数都完备，则两范数等价。
\end{exercise}
% SUPPLEMENT: id=v14-ch05-ex32-solution; parent=ch05-ex32-comparable-banach-norms; kind=solution+motivation+hypotheses; origin=independent-derivation
\begin{solution}
\textbf{补写解答（重构新增）。}恒等映射 $I:(\mathcal X,\|\cdot\|_2)\to(\mathcal X,\|\cdot\|_1)$ 为线性双射，并由给定不等式有界。两端完备，因此有界逆推论给某 $C>0$，使 $\|x\|_2\le C\|x\|_1$ 对所有 $x$ 成立。结合原界即得双边范数比较，故等价。零空间中任取 $C=1$ 即可。
\end{solution}
\begin{remark}
\textbf{动机与条件（重构新增）。}一个方向的连续性在两个 Banach 范数间自动补出反方向的连续性。可比性给恒等映射有界，两种完备性使逆定理适用；单有其中一种完备性不能用这个论证保证等价。
\end{remark}
% END SUPPLEMENT: id=v14-ch05-ex32-solution
% SOURCE p=179; printed=164; object=ch05-ex33-no-slowest-decay; number=33; region=lower
\begin{exercise}\label{exercise:ch05-no-slowest-decay}
绝对收敛级数的项不存在最慢的衰减速度：即不存在正数列 $\{a_n\}$，使 $\sum_na_n|c_n|<\infty$ 当且仅当 $\{c_n\}$ 有界。（有界数列构成 $B(\mathbb N)$，即 $\mathbb N$ 上的有界函数空间；绝对可和数列构成 $L^1(\mu)$，其中 $\mu$ 为计数测度。若存在这样的 $\{a_n\}$，考虑 $T:B(\mathbb N)\to L^1(\mu)$，定义为 $Tf(n)=a_nf(n)$。有限支集的函数集在 $L^1(\mu)$ 中稠密，却在 $B(\mathbb N)$ 中不稠密。）
\end{exercise}
% SUPPLEMENT: id=v14-ch05-ex33-solution; parent=ch05-ex33-no-slowest-decay; kind=solution+motivation+hypotheses; origin=independent-derivation
\begin{solution}
\textbf{依提示补写（重构新增）。}反设有这样的正数列。取 $c_n=1$ 得 $A=\sum_na_n<\infty$，从而 $T$ 线性且 $\|Tf\|_1\le A\|f\|_\infty$。各 $a_n>0$ 保证 $T$ 单射。对任意 $g\in\ell^1$，令 $f(n)=g(n)/a_n$，则 $\sum_na_n|f(n)|=\|g\|_1<\infty$；假设的反方向说明 $f$ 有界，因此 $T$ 满射。$B(\mathbb N)$ 与 $\ell^1$ 都完备，由有界逆定理，$T^{-1}$ 连续。

对任意 $f\in B(\mathbb N)$，$Tf$ 的有限截断 $g_N$ 在 $\ell^1$ 中趋于 $Tf$。连续逆给 $T^{-1}g_N\to f$ 于上确界范数，而各 $T^{-1}g_N$ 都有限支集。因此有限支集函数应在 $B(\mathbb N)$ 中稠密。但常数序列1与任意有限支集函数的上确界距离至少1，矛盾。
\end{solution}
\begin{remark}
\textbf{动机与条件（重构新增）。}假设把“加权可和”和“有界”完全等同，会构造出不可能存在的 Banach 空间同构，因为两空间的有限支集逼近性质不同。正权重保证可逐坐标相除；双向条件分别给 $T$ 的有界性和满射。也可直接核验矛盾：由 $\sum a_n<\infty$ 选严格递增 $n_j$ 使 $a_{n_j}<2^{-j}/j$，令 $c_{n_j}=j$、其余为0，就得到无界但加权可和的序列。
\end{remark}
% END SUPPLEMENT: id=v14-ch05-ex33-solution
% SOURCE p=179; printed=164; object=ch05-ex34-sobolev-inclusion; number=34; region=lower
\begin{exercise}\label{exercise:ch05-sobolev-inclusion}
参照\cref{exercise:ch05-ck-banach,exercise:ch05-sobolev-banach}，证明 $L_k^1([0,1])$ 到 $C^{k-1}([0,1])$ 的包含映射连续，分别使用：
\begin{enumerate}
\item\label{item:ch05-inclusion-closed-graph} 闭图像定理；
\item\label{item:ch05-inclusion-direct-bound} 直接计算。
\end{enumerate}
（本题用来说明闭图像定理可以节省工作。）
\end{exercise}
% SUPPLEMENT: id=v14-ch05-ex34-solution; parent=ch05-ex34-sobolev-inclusion; kind=solution+motivation+hypotheses; origin=independent-derivation
\begin{solution}
\textbf{两种证明（重构新增）。}依原定义 $k\ge1$ 为整数，记包含映射为 $I$。
\begin{enumerate}
\item 对应\subcref{exercise:ch05-sobolev-inclusion}{item:ch05-inclusion-closed-graph}：两空间已由所引习题证明为 Banach 空间。若 $f_n\to f$ 于 $L_k^1$ 且 $If_n\to g$ 于 $C^{k-1}$，则
\begin{align*}
\|f-g\|_1\le\|f-f_n\|_1+\|f_n-g\|_1
\le\|f-f_n\|_1+\|f_n-g\|_\infty\longrightarrow0.
\end{align*}
所以 $f=g$ 几乎处处；两者连续，故处处相等（连续非零差若存在，会在一段正长度区间上与零保持正距离）。于是 $g=If$，图像闭。闭图像定理给 $I$ 有界，从而连续。
\item 对应\subcref{exercise:ch05-sobolev-inclusion}{item:ch05-inclusion-direct-bound}：对绝对连续函数 $u$，由 $u(t)-u(s)=\int_s^tu'(v)\,\mathrm dv$ 得 $|u(t)|\le|u(s)|+\|u'\|_1$；对 $s\in[0,1]$ 积分，得 $\|u\|_\infty\le\|u\|_1+\|u'\|_1$。当 $0\le j\le k-1$，$f^{(j)}$ 均绝对连续：最高这一阶由定义，较低阶为 $C^1$ 函数，在紧区间上由微积分基本定理亦绝对连续。因此
\begin{align*}
\|If\|_{C^{k-1}}&=\sum_{j=0}^{k-1}\|f^{(j)}\|_\infty\\
&\le\sum_{j=0}^{k-1}\bigl(\|f^{(j)}\|_1+\|f^{(j+1)}\|_1\bigr)
\le2\sum_{j=0}^k\|f^{(j)}\|_1=2\|f\|_{L_k^1}.
\end{align*}
这给直接的有界估计，故连续。
\end{enumerate}
\end{solution}
\begin{remark}
\textbf{动机与条件（重构新增）。}闭图像方法只比较两个已有极限，零阶 $L^1$ 收敛加一致收敛就能辨认它们相同；具体范数界由定理自动给出。直接方法利用区间长度为1和绝对连续积分公式，给出常数2。正整数 $k$、定义中的绝对连续性以及两个已证明的完备性分别保证空间定义、直接估计和闭图像定理的适用。
\end{remark}
% END SUPPLEMENT: id=v14-ch05-ex34-solution
% SOURCE p=179; printed=164; object=ch05-ex35-closed-range-quotient; number=35; region=lower
\begin{exercise}\label{exercise:ch05-closed-range-quotient}
设 $\mathcal X,\mathcal Y$ 为 Banach 空间，$T\in L(\mathcal X,\mathcal Y)$，$\mathcal N(T)=\{x:Tx=0\}$，$\mathcal M=\operatorname{range}(T)$。则 $\mathcal X/\mathcal N(T)$ 与 $\mathcal M$ 同构，当且仅当 $\mathcal M$ 闭。（见\cref{exercise:ch05-kernel-factorization}。）
\end{exercise}
% SUPPLEMENT: id=v14-ch05-ex35-solution; parent=ch05-ex35-closed-range-quotient; kind=solution+motivation+hypotheses; origin=independent-derivation
\begin{solution}
\textbf{补写解答（重构新增）。}核为闭子空间，故 $\mathcal X/\mathcal N(T)$ 是 Banach 空间。由所引习题，$S(x+\mathcal N(T))=Tx$ 给到 $\mathcal M$ 的有界线性双射。若 $\mathcal M$ 闭，它作为 $\mathcal Y$ 的闭子空间完备，有界逆推论给 $S^{-1}$ 有界，所以 $S$ 是同构。

反过来，若存在赋范空间同构 $A:\mathcal X/\mathcal N(T)\to\mathcal M$，任意 $\mathcal M$ 中的 Cauchy 列 $y_n$ 经有界逆 $A^{-1}$ 拉回为商空间 Cauchy 列，故收敛到某 $u$；再由 $A$ 连续得 $y_n\to Au\in\mathcal M$。因此 $\mathcal M$ 完备，作为 $\mathcal Y$ 的子空间必闭。这一方向甚至不要求同构恰为自然诱导的 $S$。
\end{solution}
\begin{remark}
\textbf{动机与条件（重构新增）。}商掉核后总能得到到值域的代数双射；值域闭性使它成为两个 Banach 空间间的有界双射，从而补出逆界。$T$ 有界使核闭，$\mathcal X$ 完备使商完备，$\mathcal Y$ 完备使闭值域完备。这里“同构”沿用有界且逆有界的定义，不能换成纯代数同构。
\end{remark}
% END SUPPLEMENT: id=v14-ch05-ex35-solution
\Needspace{10\baselineskip}
% SOURCE p=179; printed=164; object=ch05-ex36-l1-universal-quotient; number=36; region=bottom
\begin{exercise}\label{exercise:ch05-l1-universal-quotient}
设 $\mathcal X$ 为可分 Banach 空间，$\mu$ 为 $\mathbb N$ 上的计数测度。设 $\{x_n\}_{n=1}^{\infty}$ 为 $\mathcal X$ 单位球的可数稠密子集，定义 $T:L^1(\mu)\to\mathcal X$ 为 $Tf=\sum_{n=1}^{\infty}f(n)x_n$。
\begin{enumerate}
\item\label{item:ch05-universal-map-bounded} $T$ 有界。
% SOURCE p=180; printed=165; object=ch05-ex36-l1-universal-quotient-continuation; region=top
\item\label{item:ch05-universal-map-onto} $T$ 满射。
\item\label{item:ch05-universal-quotient} $\mathcal X$ 同构于 $L^1(\mu)$ 的某个商空间。（使用\cref{exercise:ch05-closed-range-quotient}。）
\end{enumerate}
\end{exercise}
% SUPPLEMENT: id=v14-ch05-ex36-solution; parent=ch05-ex36-l1-universal-quotient; kind=solution+well-definedness+motivation+hypotheses; origin=independent-derivation
\begin{solution}
\textbf{逐项完整补写（重构新增）。}
\begin{enumerate}
\item 对应\subcref{exercise:ch05-l1-universal-quotient}{item:ch05-universal-map-bounded}：因 $\|x_n\|\le1$，有 $\sum_n\|f(n)x_n\|\le\sum_n|f(n)|=\|f\|_1<\infty$。$\mathcal X$ 完备保证此级数有和，所以 $T$ 良定义。有限部分和线性，取极限保持线性，且 $\|Tf\|\le\|f\|_1$，故有界。
\item 对应\subcref{exercise:ch05-l1-universal-quotient}{item:ch05-universal-map-onto}：给定 $y\in\mathcal X$，若 $y=0$ 取 $f=0$。否则令 $r_0=y$。若 $r_{j-1}\ne0$，则 $r_{j-1}/(2\|r_{j-1}\|)$ 在开单位球内，故由稠密性可选 $n_j$，使
\begin{align*}
\left\|\frac{r_{j-1}}{2\|r_{j-1}\|}-x_{n_j}\right\|<\tfrac14.
\end{align*}
令 $a_j=2\|r_{j-1}\|$、$r_j=r_{j-1}-a_jx_{n_j}$，则 $\|r_j\|<\|r_{j-1}\|/2$。若某一步残差为零，后续均令 $a_j=0$、$n_j=1$、$r_j=0$。于是对全部 $j\ge1$，
\begin{align*}
\|r_j\|\le2^{-j}\|y\|,\qquad a_j\le2^{2-j}\|y\|,\qquad \sum_ja_j\le4\|y\|.
\end{align*}
在 $\ell^1$ 中，级数 $\sum_ja_je_{n_j}$ 绝对收敛，故有和 $f\in\ell^1$；允许指标 $n_j$ 重复。对有限部分和 $f_J$ 有 $Tf_J=\sum_{j=1}^Ja_jx_{n_j}=y-r_J$。由 $T$ 有界且 $f_J\to f$，得 $Tf=\lim_JTf_J=y$，所以满射。
\item 对应\subcref{exercise:ch05-l1-universal-quotient}{item:ch05-universal-quotient}：$T$ 的值域就是闭空间 $\mathcal X$，因此前一习题给 $\ell^1/\ker T$ 与 $\mathcal X$ 同构。核闭，商范数定义良好；结论为赋范空间同构，并不在本证明中要求等距。
\end{enumerate}
\end{solution}
\begin{remark}
\textbf{动机与逐条件说明（重构新增）。}可分性提供可数个单位方向，$\ell^1$ 系数控制这些方向的绝对和。只知线性张成稠密尚不等于满射；几何衰减的残差构造让每个目标都有真正可和的原像。完备性用于定义向量级数的和。构造选半径为 $1/2$ 的归一化目标，因此单位球无论按开球还是闭球理解都适用；零空间也由零映射涵盖。
\end{remark}
% END SUPPLEMENT: id=v14-ch05-ex36-solution
% SOURCE p=180; printed=165; object=ch05-ex37-weak-tests-bounded; number=37; region=upper
\begin{exercise}\label{exercise:ch05-weak-tests-bounded}
设 $\mathcal X,\mathcal Y$ 为 Banach 空间。若线性映射 $T:\mathcal X\to\mathcal Y$ 满足每个 $f\in\mathcal Y^*$ 都有 $f\circ T\in\mathcal X^*$，则 $T$ 有界。
\end{exercise}
% SUPPLEMENT: id=v14-ch05-ex37-solution; parent=ch05-ex37-weak-tests-bounded; kind=solution+motivation+hypotheses; origin=independent-derivation
\begin{solution}
\textbf{补写解答（重构新增）。}设 $x_n\to x$ 且 $Tx_n\to y$。对每个 $f\in\mathcal Y^*$，连续性及题设给
\begin{align*}
f(y)=\lim_nf(Tx_n)=\lim_n(f\circ T)(x_n)=(f\circ T)(x)=f(Tx).
\end{align*}
对偶分离点，所以 $y=Tx$。故 $T$ 图像闭，由闭图像定理得到有界。

也可用一致有界原理直接看出目标完备性可省。考虑 $\mathcal A=\{f\circ T:f\in\mathcal Y^*,\ \|f\|\le1\}\subseteq\mathcal X^*$。对固定 $x$，有 $\sup_{A\in\mathcal A}|Ax|\le\|Tx\|<\infty$。$\mathcal X$ 完备给 $\sup_{A\in\mathcal A}\|A\|=C<\infty$。范数支撑泛函给
\begin{align*}
\|Tx\|=\sup_{\|f\|\le1}|f(Tx)|\le C\|x\|,
\end{align*}
所以 $T$ 有界。
\end{solution}
\begin{remark}
\textbf{动机与条件（重构新增）。}每个连续标量测量都认可正确的极限，而这些测量能分离向量，故图像闭。第一种证明按题设使用两端完备性；第二种证明表明实际只需定义域完备，目标赋范即可。题设并没有给所有 $f\circ T$ 一个共同的界，这个共同界由一致有界原理产生。
\end{remark}
% END SUPPLEMENT: id=v14-ch05-ex37-solution
% SOURCE p=180; printed=165; object=ch05-ex38-pointwise-operator-limit; number=38; region=upper
\begin{exercise}\label{exercise:ch05-pointwise-operator-limit}
设 $\mathcal X,\mathcal Y$ 为 Banach 空间，$\{T_n\}\subseteq L(\mathcal X,\mathcal Y)$，且对每个 $x\in\mathcal X$，$\lim_nT_nx$ 都存在。令 $Tx=\lim_nT_nx$，则 $T\in L(\mathcal X,\mathcal Y)$。
\end{exercise}
% SUPPLEMENT: id=v14-ch05-ex38-solution; parent=ch05-ex38-pointwise-operator-limit; kind=solution+motivation+hypotheses; origin=independent-derivation
\begin{solution}
\textbf{补写解答（重构新增）。}对固定 $x$，收敛的序列 $T_nx$ 有界，所以 $\sup_n\|T_nx\|<\infty$。一致有界原理给某 $M<\infty$，使每个 $n$ 都有 $\|T_n\|\le M$。由加法和数乘的连续性，
\begin{align*}
T(ax+by)=\lim_nT_n(ax+by)=a\lim_nT_nx+b\lim_nT_ny=aTx+bTy,
\end{align*}
所以 $T$ 线性。范数连续性再给 $\|Tx\|=\lim_n\|T_nx\|\le M\|x\|$，故有界。
\end{solution}
\begin{remark}
\textbf{动机与条件（重构新增）。}逐点极限虽然没有直接给统一算子界，定义域的完备性使一致有界原理补出这个界，因而极限仍连续。既然题设已保证极限在 $\mathcal Y$ 中存在，本证明不再需要 $\mathcal Y$ 完备；若只假设各 $T_nx$ 为 Cauchy 列，目标完备性才负责极限存在。结论不声称 $T_n$ 按算子范数趋于 $T$。
\end{remark}
% END SUPPLEMENT: id=v14-ch05-ex38-solution
% SOURCE p=180; printed=165; object=ch05-ex39-bilinear-continuity; number=39; region=middle
\begin{exercise}\label{exercise:ch05-bilinear-continuity}
设 $\mathcal X,\mathcal Y,\mathcal Z$ 为 Banach 空间，$B:\mathcal X\times\mathcal Y\to\mathcal Z$ 为分别连续的双线性映射，即每个 $x\in\mathcal X$ 都有 $B(x,\cdot)\in L(\mathcal Y,\mathcal Z)$，每个 $y\in\mathcal Y$ 都有 $B(\cdot,y)\in L(\mathcal X,\mathcal Z)$。则 $B$ 联合连续，即作为 $\mathcal X\times\mathcal Y\to\mathcal Z$ 的映射连续。（把问题归结为证明存在 $C>0$，使 $\|B(x,y)\|\le C\|x\|\|y\|$。）
\end{exercise}
% SUPPLEMENT: id=v14-ch05-ex39-solution; parent=ch05-ex39-bilinear-continuity; kind=solution+motivation+hypotheses; origin=independent-derivation
\begin{solution}
\textbf{补写解答（重构新增）。}对 $\|x\|\le1$，记 $B_x(y)=B(x,y)$，它属于 $L(\mathcal Y,\mathcal Z)$。固定 $y$，另一个截面 $C_y(x)=B(x,y)$ 有界，所以
\begin{align*}
\sup_{\|x\|\le1}\|B_xy\|\le\|C_y\|<\infty.
\end{align*}
在完备定义域 $\mathcal Y$ 上应用一致有界原理，得 $\sup_{\|x\|\le1}\|B_x\|=M<\infty$。若 $x\ne0$，对 $x/\|x\|$ 应用此界并用双线性，得 $\|B(x,y)\|\le M\|x\|\|y\|$；$x=0$ 时同样成立。可取 $C=M+1>0$。

在任意 $(x_0,y_0)$ 附近，
\begin{align*}
\|B(x,y)-B(x_0,y_0)\|
&\le\|B(x-x_0,y)\|+\|B(x_0,y-y_0)\|\\
&\le C\bigl(\|x-x_0\|\|y\|+\|x_0\|\|y-y_0\|\bigr).
\end{align*}
当 $\|y-y_0\|<1$ 时，$\|y\|\le\|y_0\|+1$，于是右端在 $(x,y)\to(x_0,y_0)$ 时趋零，证明联合连续。
\end{solution}
\begin{remark}
\textbf{动机与逐条件说明（重构新增）。}一个方向的连续性提供有界算子族，另一个方向的连续性给这族的逐点界；一致有界原理把两者组合成双线性的统一乘积界。按此证明只需 $\mathcal Y$ 完备，$\mathcal X,\mathcal Z$ 赋范即可；交换变量也可改用 $\mathcal X$ 完备。双线性用于归一化和两项误差分解，一般分别连续函数没有同样结论。
\end{remark}
% END SUPPLEMENT: id=v14-ch05-ex39-solution
% SOURCE p=180; printed=165; object=ch05-ex40-condensation-singularities; number=40; region=middle
\begin{exercise}[奇性凝聚原理]\label{exercise:ch05-condensation-singularities}
设 $\mathcal X,\mathcal Y$ 为 Banach 空间，$\{T_{jk}:j,k\in\mathbb N\}\subseteq L(\mathcal X,\mathcal Y)$。若对每个 $k$ 都存在 $x\in\mathcal X$，使 $\sup_{j\in\mathbb N}\|T_{jk}x\|=\infty$，则存在某个 $x$（事实上，存在一个剩余的这样的 $x$ 的集合），使所有 $k$ 都满足 $\sup_{j\in\mathbb N}\|T_{jk}x\|=\infty$。
\end{exercise}
% SUPPLEMENT: id=v14-ch05-ex40-solution; parent=ch05-ex40-condensation-singularities; kind=solution+motivation+hypotheses; origin=independent-derivation
\begin{solution}
\textbf{补写解答（重构新增）。}对每个固定 $k$，令 $S_k=\{x:\sup_j\|T_{jk}x\|<\infty\}$。若 $S_k$ 非疏，则一致有界原理第一问给 $M_k=\sup_j\|T_{jk}\|<\infty$，从而每个 $x$ 都有 $\sup_j\|T_{jk}x\|\le M_k\|x\|<\infty$，与题设存在一个发散点矛盾。因此每个 $S_k$ 疏，可数并 $\bigcup_kS_k$ 也疏。其补
\begin{align*}
R=\bigcap_{k=1}^{\infty}\{x:\sup_j\|T_{jk}x\|=\infty\}
\end{align*}
是剩余集。$\mathcal X$ 为非空完备度量空间，故剩余集稠密且非空；每个 $x\in R$ 同时满足全部所需发散条件。
\end{solution}
\begin{remark}
\textbf{动机与条件（重构新增）。}每一列算子只要在某一点失控，就迫使其逐点有界集拓扑上很小；可数多个这样的障碍仍可同时避开。指标 $k$ 可数用于疏集的可数并；定义域完备性使剩余集非空。目标的完备性在本证明中未使用，因为只应用一致有界原理并考察范数，不在目标中取新的极限。
\end{remark}
% END SUPPLEMENT: id=v14-ch05-ex40-solution
% SOURCE p=180; printed=165; object=ch05-ex41-countable-hamel-dimension; number=41; region=middle-lower
\begin{exercise}\label{exercise:ch05-countable-hamel-dimension}
设 $\mathcal X$ 为可数无限维向量空间，即每个元素是某个可数无限线性无关集中的有限多个元素的线性组合。不存在使 $\mathcal X$ 完备的范数。（给定 $\mathcal X$ 上任意范数，应用\subcref{exercise:ch05-finite-subspace-closed}{item:ch05-finite-dim-closed}及 Baire 纲定理。）
\end{exercise}
% SUPPLEMENT: id=v14-ch05-ex41-solution; parent=ch05-ex41-countable-hamel-dimension; kind=solution+motivation+hypotheses; origin=independent-derivation
\begin{solution}
\textbf{补写解答（重构新增）。}设 $\{e_n:n\ge1\}$ 为所给 Hamel 基，赋予任意范数。令 $\mathcal M_n=\operatorname{span}\{e_1,\ldots,e_n\}$，则每个 $\mathcal M_n$ 有限维，故闭；基无限且线性无关，保证它是真子空间。

真子空间没有内部：若某个 $B(x,r)\subseteq\mathcal M_n$，则 $x\in\mathcal M_n$，平移得 $B(0,r)\subseteq\mathcal M_n$。对任意 $y\ne0$，选 $t>0$ 足够小使 $\|ty\|<r$，就有 $ty\in\mathcal M_n$，再除以 $t$ 得 $y\in\mathcal M_n$，矛盾。因此各 $\mathcal M_n$ 闭且无处稠密。

Hamel 展开只用有限多个基向量，所以 $\mathcal X=\bigcup_{n\ge1}\mathcal M_n$。若该范数完备，这与非空完备空间的 Baire 纲定理矛盾。因此没有这样的完备范数。
\end{solution}
\begin{remark}
\textbf{动机与条件（重构新增）。}可数无限的代数基会把整个空间拆成可数个有限维闭薄层，Baire 定理不允许完备空间这样被覆盖。这里的“可数维”指有限线性组合意义的 Hamel 维数，不能与允许无限收敛展开的可数 Schauder 基混同；后者可以出现在无限维 Banach 空间中。
\end{remark}
% END SUPPLEMENT: id=v14-ch05-ex41-solution
\Needspace{12\baselineskip}
% SOURCE p=180; printed=165; object=ch05-ex42-nowhere-differentiable; number=42; region=lower
\begin{exercise}\label{exercise:ch05-nowhere-differentiable}
\begin{definition*}[本题的差商有界集合]\phantomsection\label{definition:ch05-difference-quotient-sets}
令 $E_n$ 为所有满足如下性质的 $f\in C([0,1])$ 所成集合：存在一个依赖于 $f$ 的 $x_0\in[0,1]$，使所有 $x\in[0,1]$ 都满足 $|f(x)-f(x_0)|\le n|x-x_0|$。
\end{definition*}
\begin{enumerate}
\item\label{item:ch05-bounded-quotient-nowhere-dense} $E_n$ 在 $C([0,1])$ 中无处稠密。（任意实值 $f\in C([0,1])$ 可被分段线性函数 $g$ 一致逼近，$g$ 有有限多个线性段且斜率为 $\pm2n$。若 $\|h-g\|_u$ 足够小，则 $h\notin E_n$。）
\item\label{item:ch05-nowhere-differentiable-residual} 处处不可微函数构成 $C([0,1])$ 中的剩余集。
\end{enumerate}
\end{exercise}
% SUPPLEMENT: id=v14-ch05-ex42-solution; parent=ch05-ex42-nowhere-differentiable; kind=solution+source-correction+counterexample+motivation+hypotheses; origin=independent-derivation
\begin{remark}
\textbf{原提示中的斜率修正（重构新增）。}原页明确写固定斜率 $\pm2n$；按字面这不足以逼近任意连续函数。所有这种有限折线都是 $2n$-Lipschitz 的。若它与 $f(t)=3nt$ 的一致误差为 $\varepsilon$，比较两个端点就有 $3n\le2n+2\varepsilon$，即 $\varepsilon\ge n/2$，不能任意小。可行的修正是允许选择依赖于目标逼近的 $M>2n$，让各线性段斜率为 $\pm M$。下面完整构造这种折线；本题 $n$ 取正整数。
\end{remark}
\begin{solution}
\textbf{逐项完整补写（重构新增）。}
\begin{enumerate}
\item 对应\subcref{exercise:ch05-nowhere-differentiable}{item:ch05-bounded-quotient-nowhere-dense}：先证 $E_n$ 闭。设 $f_j\in E_n$ 且 $f_j\to f$ 一致，取相应见证点 $t_j$。由 $[0,1]$ 紧，选子列使 $t_j\to t_0$，仍沿用原指标。因
\begin{align*}
|f_j(t_j)-f(t_0)|\le\|f_j-f\|_u+|f(t_j)-f(t_0)|\longrightarrow0,
\end{align*}
对任意固定 $x$，在 $|f_j(x)-f_j(t_j)|\le n|x-t_j|$ 中取极限，得 $|f(x)-f(t_0)|\le n|x-t_0|$。同一个 $t_0$ 适用于全部 $x$，故 $f\in E_n$。

再证它没有内部，且给出提示所需的稳健逼近。给任意 $f\in C([0,1])$ 和误差 $\eta>0$，写 $f=u+iv$，实值情形令 $v=0$。由 $u$ 一致连续，细分区间并在端点线性插值，得到实折线 $p$，使 $\|p-u\|_u<\eta/4$：在每个小区间上，插值是两端值的凸组合，两端值与该点的 $u$ 值都相差小于 $\eta/4$ 即可。

令 $L$ 为 $p$ 各有限线性段斜率绝对值的最大值，取 $M>\max(L,2n)$。将每段再有限细分成长度 $h<\eta/(16M)$ 的小段。若某小段的原斜率为 $a$，取 $\theta=(1+a/M)/2\in(0,1)$，用前 $\theta h$ 长度上斜率 $+M$、后 $(1-\theta)h$ 长度上斜率 $-M$ 的两段折线替换；从原左端值起，净增量为 $M\theta h-M(1-\theta)h=ah$，所以仍到达原右端值。这样拼接成有限折线 $q$，各段斜率为 $\pm M$。每个小段中 $|q-p|\le(M+L)h<\eta/8$，故 $\|q-u\|_u<3\eta/8<\eta/2$。

令 $g=q+iv$，则 $\|g-f\|_u<\eta/2$。$q$ 有有限多个正长度线性段，记最小段长为 $\ell>0$。对任何 $t\in[0,1]$，选一个包含它的线性段，再选离它较远的端点 $s$，则 $|s-t|\ge\ell/2$，且 $|q(s)-q(t)|=M|s-t|$。取
\begin{align*}
0<\rho\le\min\bigl(\eta/2,(M-n)\ell/8\bigr).
\end{align*}
若 $\|h-g\|_u<\rho$，则利用实部估计得到
\begin{align*}
|h(s)-h(t)|&\ge|\operatorname{Re}h(s)-\operatorname{Re}h(t)|\\
&\ge M|s-t|-2\rho>n|s-t|.
\end{align*}
最后严格不等式来自 $2\rho\le(M-n)\ell/4<(M-n)|s-t|$。所以每个 $t$ 都不能作为 $E_n$ 的见证点，$h\notin E_n$。同时 $\|h-f\|_u<\eta$，因此任意 $f$ 的任意球内都含一个与 $E_n$ 不交的小球。结合闭性，$E_n$ 无处稠密。
\item 对应\subcref{exercise:ch05-nowhere-differentiable}{item:ch05-nowhere-differentiable-residual}：若 $f$ 在某 $t_0$ 具有有限导数（端点用区间内的单侧导数），则存在 $C<\infty$、$\delta>0$，使 $0<|t-t_0|<\delta$ 时 $|f(t)-f(t_0)|\le C|t-t_0|$。而 $|t-t_0|\ge\delta$ 时，有
\begin{align*}
|f(t)-f(t_0)|\le2\|f\|_u\le\frac{2\|f\|_u}{\delta}|t-t_0|.
\end{align*}
选正整数 $n\ge\max(C,2\|f\|_u/\delta)$，就有 $f\in E_n$。因此至少一点可微的函数集包含于疏集 $\bigcup_{n\ge1}E_n$，所以它本身疏。其补集，即处处不可微的函数集，为剩余集。$C([0,1])$ 完备且非空，故此剩余集稠密并非空，确实给出这类函数的存在性。
\end{enumerate}
\end{solution}
\begin{remark}
\textbf{动机与逐条件说明（重构新增）。}“在某点有导数”先转成该点差商的某个整数界；可数个整数界使坏集合可被可数个闭薄层覆盖。紧区间使见证点可取收敛子列，并使局部导数界在远离该点处借函数有界性扩成全区间界。折线有限分段给最小正段长，因而能选一个同时排除所有候选点的误差半径，不能只对每个点分别选半径。复值情形只对实部加折线扰动，保留虚部，同样得到模长差商下界。第二问只用“有一点可微蕴含属于某个 $E_n$”，并未误称这两个条件等价。
\end{remark}
% END SUPPLEMENT: id=v14-ch05-ex42-solution
\end{document}
